\chapter{Abstract}

Educational materials at HCMUT—including lecture videos, audio recordings, slides, and documents—are expanding rapidly, yet conventional text-based tools fail to capture their multimodal nature. These methods often rely on transcripts or extracted text, resulting in recognition errors, hallucinations, and loss of visual and auditory context.

This project proposes a web-based application powered by a Multimodal Retrieval-Augmented Generation (RAG) pipeline to assist students in studying and exam preparation. The system integrates textual transcripts, OCR-based slide and document extraction, visual information for diagrams and equations, and audio features that capture emphasis and speaker prosody. Retrieval combines sparse and dense methods with cross-modal verification, and retrieved evidence is grounded in a large language model for accurate and explainable responses.

The prototype includes a backend using [PLACEHOLDER] for multimodal vector storage and a frontend built with [PLACEHOLDER] for responsive search and summarization. Key functionalities include multimodal content retrieval, automatic lecture summarization, and personalized study roadmap generation. The system aims to enhance learning efficiency, reduce information loss, and provide a robust, adaptive educational assistant tailored to the HCMUT academic environment.

\newpage
